\documentclass[11pt,a4paper]{article}
\usepackage[T1]{fontenc}
\usepackage[utf8]{inputenc}
\usepackage{amsmath,amssymb,xcolor}
\usepackage[margin=25mm]{geometry}
\usepackage[colorlinks=true,urlcolor=blue]{hyperref}
\definecolor{nexorablue}{HTML}{2459B5}
\setlength{\parindent}{0pt}
\setlength{\parskip}{7pt}
\setlength{\emergencystretch}{2em}
\renewcommand{\familydefault}{\sfdefault}
\begin{document}

{\small\bfseries\color{nexorablue} NEXORA CPGE}\par
{\small MP / NEX-TOP-04 / Énoncé}\par
{\LARGE\bfseries Une majoration matricielle atteinte}\par
\vspace{0.25em}\hrule\vspace{0.6em}

Munir $\mathbb R^3$ et $\mathbb R^2$ de la norme maximum. Calculer la norme de l’opérateur $u(x,y,z)=(2x-y+3z,-x+4y+z)$ et exhiber un vecteur unité sur lequel elle est atteinte. En déduire une constante de Lipschitz pour $u$.


\end{document}
